\begin{center}
\Large \textbf{Les sylphes}
\end{center}

Un formulaire et des données numériques sont fournis à la fin du sujet.
Toutes les applications numériques seront effectuées avec 1 chiffre significatif.

Les sylphes sont des émissions lumineuses rouges, très intenses et relativement brèves (de l'ordre de quelques centaines de millisecondes) qui se produisent entre le sommet des nuages et l'ionosphère. La majorité d'entre-elles se situe entre 40 et 80 km d'altitude. Elles ne sont observées et étudiées que depuis 1990 environ. Certaines de leurs manifestations avaient déjà été rapportées bien avant mais leur observation depuis la Terre est difficile. Les sylphes font partie d'un groupe d'émissions lumineuses qui se produisent au moment des orages. On retrouve dans le même groupe d'émissions les jets bleus et les elfes, voir la figure 1. Ces émissions sont la conséquence de courants électriques importants qui se produisent lors d'une avalanche de créations d'ions et d'électrons. L'initiateur de l'avalanche peut être une particule du rayonnement cosmique. Lors de chocs avec les molécules, les électrons des courants électriques apportent de l'énergie à celles-ci, perdue ensuite en émettant de la lumière visible.

\TLEFigure{1}{Différentes formes de phénomènes lumineux éphémères atmosphériques}

Ce sujet comporte 3 parties largement indépendantes. La première concerne l'observation des sylphes. La seconde étudie des propriétés électriques de l'atmosphère. La dernière partie présente le modèle de la formation avalancheuse du courant à l'origine des sylphes.

\section*{I Observation des sylphes}

\begin{itemize}
    \item[$\square$ -- 1.] À l'aide des courbes de la figure 2, justifier la couleur caractéristiques des sylphes. Justifier également leur observation difficile depuis la surface de la Terre.
\end{itemize}

Malgré cette difficulté, des images de sylphes ont été obtenues comme celles de la figure 3 qui ont été observées depuis l'Observatoire du Pic du Midi de Bigorre. Des sylphes y ont été vues au-dessus des Alpes. Le Pic du Midi de Bigorre se situe au milieu de la chaîne des Pyrénées qui s'étend de Perpignan à l'Est jusqu'à Biarritz à l'Ouest, séparant la France de l'Espagne.
